% Part: many-valued-logic
% Chapter: three-valued-logics
% Section: lukasiewicz

\documentclass[../../../include/open-logic-section]{subfiles}

\begin{document}

\olfileid{mvl}{inf}{god}

\olsection{G\"odel तर्कशास्त्रे}

\begin{editorial}
  अनंतमूल्य G\"odel तर्कशास्त्रावरील हा विभाग
  संक्षिप्त ``प्रारूप'' आहे.
\end{editorial}

\begin{defn}\ollabel{def:goedel} अनंतमूल्य G\"odel
तर्कशास्त्र~$\LogGod[\infty]$ पुढील मॅट्रिक्सने परिभाषित होते:
\begin{enumerate}
  \item $\lfalse$, $\lnot$, $\land$, $\lor$, $\lif$ असलेली मानक विधानीय भाषा $\Lang L_0$.
  \item सत्यतामूल्यांचा संच $V_\infty$.
  \item $1$ हेच एकमेव निर्दिष्ट मूल्य आहे, म्हणजे $V^+ = \{1\}$.
  \item सत्यता-फलने पुढील सूत्रांनी दिली आहेत:
  \begin{align*}
    \tf{\lfalse} & = 0\\
    \tf{\lnot}[\LogGod](x) & = \begin{cases}
      1 & \text{जर } x =0\\
      0 & \text{इतर वेळी}
    \end{cases}\\
    \tf{\land}[\LogGod](x,y) & = \min(x,y)\\
    \tf{\lor}[\LogGod](x,y) & = \max(x,y)\\
    \tf{\lif}[\LogGod](x,y) & = \begin{cases}
      1 & \text{जर } x \le y\\
      y & \text{इतर वेळी.}
    \end{cases}
    \end{align*}
\end{enumerate}
$m$-मूल्य G\"odel तर्कशास्त्राची व्याख्या याचप्रमाणे,
फक्त $V = V_m$ घेऊन केली जाते.
\end{defn}

\begin{prop}
  \olref[thr][god]{defn:goedel} मध्ये परिभाषित
  $\LogGod[3]$ आणि \olref{def:goedel} मध्ये परिभाषित
  $\LogGod[3]$ हे एकच तर्कशास्त्र आहे.
\end{prop}

\begin{proof}
  \olref[thr][god]{defn:goedel} मधील संयोजकांची
  सत्यताकोष्टके आणि \olref{def:goedel} मधील
  समीकरणांनी निश्चित होणारी सत्यताकोष्टके यांची तुलना
  केल्यास हे दिसते:
  \begin{center}
    \begin{tabular}{c|c}
      $\tf{\lnot}[\LogGod[3]]$ & \\
      \hline
      $1$ & $0$ \\
      $1/2$ & $0$ \\
      $0$ & $1$
    \end{tabular}
    \quad
    \begin{tabular}{c|ccc}
      $\tf{\land}[\LogGod]$ & $1$ & $1/2$ & $0$ \\
      \hline
      $1$ & $1$ & $1/2$ & $0$ \\
      $1/2$ & $1/2$ & $1/2$ & $0$\\
      $0$ & $0$ & $0$ & $0$
    \end{tabular}
    \\[2ex]
    \begin{tabular}{c|ccc}
      $\tf{\lor}[\LogGod]$ & $1$ & $1/2$ & $0$ \\
      \hline
      $1$ & $1$ & $1$ & $1$ \\
      $1/2$ & $1$ & $1/2$ & $1/2$ \\
      $0$ & $1$ & $1/2$ & $0$
    \end{tabular}
    \quad
    \begin{tabular}{c|ccc}
      $\tf{\lif}[\LogGod]$ & $1$ & $1/2$ & $0$ \\
      \hline
      $1$ & $1$ & $1/2$ & $0$ \\
      $1/2$ & $1$ & $1$ & $0$  \\
      $0$ & $1$ & $1$ & $1$
    \end{tabular}
  \end{center}
\end{proof}

\begin{prop}\ollabel{prop:god-infty-m}
  $\Gamma \Entails[\LogGod[\infty]] !B$ असेल, तर
  सर्व~$m \ge 2$ साठी $\Gamma \Entails[\LogGod[m]] !B$.
\end{prop}

\begin{proof}
  सराव.
\end{proof}

\begin{prob}
  \olref[mvl][inf][god]{prop:god-infty-m} सिद्ध करा.
\end{prob}

खरे तर आधारसूत्रांचा संच ससीम असेल, तर याचे प्रतिलोम
विधानही खरे आहे.

$\LogGod[3]$ प्रमाणेच $\LogGod[\infty]$ मध्येही
अंतःप्रज्ञावादी तर्कशास्त्रात वैध असलेली सर्व !!{formula}s
सर्वतः सत्य आहेत; आणि आधीची सर्वतः सत्य नसलेल्या सूत्रांची
उदाहरणे $\LogGod[\infty]$ मध्येही सर्वतः सत्य नाहीत:
\begin{align*}
  & p \lor \lnot p && (p \lif q) \lif (\lnot p \lor q) \\
  & \lnot\lnot p \lif p && \lnot(\lnot p \land \lnot q) \lif (p \lor q) \\
  & ((p \lif q) \lif p) \lif p && \lnot(p \lif q) \lif (p \land \lnot q)
\end{align*}
अंतःप्रज्ञावादी तर्कशास्त्रात वैध नसलेले, पण
$\LogGod[3]$ मध्ये सर्वतः सत्य असलेले
!!{formula}~$(p \lif q) \lor (q \lif p)$ हे
$\LogGod[\infty]$ मध्येही सर्वतः सत्य आहे. खरे तर
$(!A \lif !B) \lor (!B \lif !A)$ ही योजना
अंतःप्रज्ञावादी तर्कशास्त्रात जोडून $\LogGod[\infty]$
चे वैशिष्ट्यीकरण करता येते. हे Michael Dummett यांनी
दाखवले; म्हणून $\LogGod[\infty]$ ला अनेकदा
G\"odel--Dummett तर्कशास्त्र~$\Log{LC}$ म्हणतात.

\begin{prob}
  $(p \lif q) \lor (q \lif p)$ हे $\LogGod[\infty]$
  मध्ये सर्वतः सत्य सूत्र आहे हे दाखवा.
\end{prob}

\begin{prob}
  $(p \lif q) \lor (q \lif r) \lor (r \lif s)$ हे
  $\LogGod[3]$ मध्ये सर्वतः सत्य असूनही
  $\LogGod[\infty]$ मध्ये सर्वतः सत्य नाही हे दाखवा.
\end{prob}

\end{document}
